\documentclass[a4paper,11pt]{article}

% ============================================================
%  Le Collé — Algèbre linéaire (compléments de deuxième année)
%  © 2026 Ahmed Amine Ouqlani
% ============================================================

% ---------- Polices ----------
\usepackage[utf8]{inputenc}
\usepackage[T1]{fontenc}
\usepackage[sc]{mathpazo}          % texte et mathématiques : Palatino
\usepackage[scaled=0.88]{helvet}   % titres et repères : linéale
\linespread{1.07}
\usepackage{amsmath,amssymb,amsthm,mathrsfs,stmaryrd}
\usepackage{microtype}

% ---------- Mise en page ----------
\usepackage[a4paper,top=2.3cm,bottom=2.3cm,left=2.4cm,right=2.4cm,
            headheight=14pt,headsep=0.8cm,footskip=1.1cm]{geometry}
\setlength{\parindent}{0pt}
\setlength{\parskip}{0.45em plus 0.15em minus 0.1em}
\raggedbottom
\emergencystretch=1.5em
\clubpenalty=10000 \widowpenalty=10000 \displaywidowpenalty=10000 \brokenpenalty=10000

% ---------- Couleurs : un bleu de nuit, un ocre, des gris ----------
\usepackage{xcolor}
\definecolor{Ink}{HTML}{1C2533}
\definecolor{Bleu}{HTML}{1F3A5F}
\definecolor{BleuMoyen}{HTML}{4A6A8F}
\definecolor{BleuClair}{HTML}{9DB2C9}
\definecolor{Fond}{HTML}{F2F4F7}
\definecolor{Ocre}{HTML}{B07A26}
\definecolor{Brique}{HTML}{A23B2A}
\definecolor{Gris}{HTML}{5E6672}
\definecolor{Filet}{HTML}{C5CCD6}

% ---------- Packages ----------
\usepackage{enumitem}
\setlist{itemsep=0.25em,topsep=0.35em,parsep=0pt}
\setlist[itemize]{label=\textcolor{BleuMoyen}{\small\textbullet},leftmargin=1.4em}
\setlist[enumerate]{leftmargin=1.8em,label=\textsf{\small\color{Bleu}\arabic*.}}
\usepackage{booktabs,array,multicol}
\usepackage{needspace,float}
\usepackage[font=small,labelfont={sf,color=Bleu},labelsep=period,skip=6pt]{caption}
\usepackage{tikz}
\usetikzlibrary{arrows.meta,calc,decorations.markings}
\usepackage{pgfplots}
\pgfplotsset{compat=1.17}
\usepackage{tkz-tab}
\usepackage[most]{tcolorbox}
\usepackage{titlesec,titletoc}
\usepackage{fancyhdr,lastpage,etoolbox}
\usepackage[colorlinks=true,linkcolor=Bleu,urlcolor=Bleu,citecolor=Bleu,
            pdftitle={Entraînement séries — samedi 28 novembre 2026},pdfauthor={Ahmed Amine Ouqlani}]{hyperref}

\numberwithin{equation}{section}
\numberwithin{figure}{section}
\numberwithin{table}{section}

% ---------- Titres ----------
% numéro du chapitre dans le parcours PT (affiché), distinct de la numérotation interne
\newcommand{\grandnum}{\thesection}\newcommand{\ongletcouleur}{Bleu}
% Chapitres (sections) : grand numéro, titre en linéale, filet ocre
\titleformat{\section}[display]
  {\normalfont\sffamily\color{Ink}}
  {{\small\color{Gris}\textls[200]{\MakeUppercase{Chapitre}}}\\[-0.2em]%
   {\fontsize{64}{64}\selectfont\color{BleuClair}\grandnum}}
  {0.4em}
  {\fontsize{25}{30}\selectfont\bfseries}
  [\vspace{0.5em}{\color{Ocre}\rule{3.2cm}{1.4pt}}]
\titlespacing*{\section}{0pt}{0pt}{1.6em}
\titleformat{name=\section,numberless}[display]
  {\normalfont\sffamily\color{Ink}}{}{0pt}{\fontsize{25}{30}\selectfont\bfseries}
  [\vspace{0.5em}{\color{Ocre}\rule{3.2cm}{1.4pt}}]
\titlespacing*{name=\section,numberless}{0pt}{0pt}{1.6em}

\titleformat{\subsection}
  {\needspace{5\baselineskip}\normalfont\sffamily\Large\bfseries\color{Ink}}
  {\color{Bleu}\thesubsection}{0.8em}{}
\titlespacing*{\subsection}{0pt}{1.8em}{0.5em}
\titleformat{\subsubsection}
  {\needspace{5\baselineskip}\normalfont\sffamily\large\bfseries\color{Ink}}
  {\color{Bleu}\thesubsubsection}{0.7em}{}
\titlespacing*{\subsubsection}{0pt}{1.6em}{0.4em}
\titleformat{\paragraph}[runin]{\normalfont\sffamily\bfseries\color{Ink}}{}{0pt}{}[.]
\titlespacing*{\paragraph}{0pt}{0.9em}{0.6em}
\setcounter{secnumdepth}{3}
\setcounter{tocdepth}{2}

% ---------- Table des matières ----------
\titlecontents{section}[0pt]{\addvspace{0.7em}\sffamily\bfseries\color{Ink}}
  {}{}{\hfill\contentspage}[\addvspace{0.2em}]
\titlecontents{subsection}[2.2em]{\small\linespread{1}\selectfont}
  {\makebox[2.6em][l]{\color{Gris}\thecontentslabel}}{}{\titlerule*[0.7pc]{.}\contentspage}

% ---------- En-têtes et pieds de page ----------
\pagestyle{fancy}
\fancyhf{}
\renewcommand{\sectionmark}[1]{\markboth{#1}{}}
\fancyhead[L]{\footnotesize\sffamily\color{Gris}\leftmark}
\fancyhead[R]{\footnotesize\sffamily\scshape\color{Gris}Le Collé}
\fancyfoot[C]{\small\sffamily\color{Ink}\thepage}
\fancyfoot[L]{\scriptsize\sffamily\hyperlink{toc}{\color{Filet}sommaire}}
\renewcommand{\headrulewidth}{0.4pt}
\renewcommand{\headrule}{{\color{Filet}\hrule height\headrulewidth width\headwidth}}
\renewcommand{\footrulewidth}{0pt}
\fancypagestyle{plain}{\fancyhf{}\fancyfoot[C]{\small\sffamily\color{Ink}\thepage}\renewcommand{\headrulewidth}{0pt}}

% ---------- Énoncés ----------
\newtheoremstyle{enonce}{0pt}{0pt}{}{}{\sffamily\bfseries\color{Bleu}}{.}{0.5em}
  {\thmname{#1}\thmnumber{~#2}\thmnote{\ \normalfont\sffamily\color{Ink}(#3)}}
\newtheoremstyle{exemple}{0.9em}{0.9em}{}{}{\sffamily\bfseries\color{BleuMoyen}}{.}{0.5em}
  {\thmname{#1}\thmnumber{~#2}\thmnote{\ \normalfont\sffamily\color{Ink}(#3)}}
\newtheoremstyle{remarque}{0.7em}{0.7em}{}{}{\sffamily\itshape}{.}{0.5em}{\thmname{#1}\thmnote{\ (#3)}}
\theoremstyle{enonce}
\newtheorem{definition}{Définition}[section]
\newtheorem{theoreme}[definition]{Théorème}
\newtheorem{proposition}[definition]{Proposition}
\newtheorem{corollaire}[definition]{Corollaire}
\theoremstyle{exemple}
\newtheorem{exemple}[definition]{Exemple}
\theoremstyle{remarque}
\newtheorem*{remarque}{Remarque}

\tcbset{
  defi/.style={enhanced,breakable,colback=white,colframe=white,boxrule=0pt,arc=0pt,
     borderline west={1.2pt}{0pt}{Bleu},left=11pt,right=4pt,top=3pt,bottom=3pt,
     before skip=1.1em,after skip=1.1em},
  thm/.style={enhanced,breakable,colback=Fond,colframe=Fond,boxrule=0pt,arc=0pt,
     borderline west={2.4pt}{0pt}{Bleu},left=11pt,right=8pt,top=6pt,bottom=6pt,
     before skip=1.2em,after skip=1.2em}}
\tcolorboxenvironment{definition}{defi}
\tcolorboxenvironment{theoreme}{thm}
\tcolorboxenvironment{proposition}{thm}
\tcolorboxenvironment{corollaire}{thm}

\renewcommand{\proofname}{Démonstration}
\renewcommand{\qedsymbol}{\textcolor{BleuClair}{$\blacksquare$}}
\makeatletter
\renewenvironment{proof}[1][\proofname]{\par\pushQED{\qed}\normalfont
  \topsep0.5em\relax\trivlist\item[\hskip\labelsep\sffamily\itshape\color{Gris}#1.]\ignorespaces}
  {\popQED\endtrivlist\@endpefalse}
\makeatother

% ---------- Méthode : filets haut et bas, titre en bandeau ----------
\makeatletter\newcommand{\methlabelfix}{\protected@edef\@currentlabel{\thetcbcounter}}\makeatother
\newtcolorbox[auto counter,number within=section]{methode}[2][]{%
  enhanced,unbreakable,colback=white,colframe=white,boxrule=0pt,arc=0pt,
  borderline north={0.9pt}{0pt}{Bleu},borderline south={0.9pt}{0pt}{Bleu},
  left=4pt,right=4pt,top=8pt,bottom=8pt,before skip=1.4em,after skip=1.4em,
  before upper={\methlabelfix\phantomsection\addcontentsline{met}{metentry}{\protect\makebox[2.6em][l]{\thetcbcounter}#2}%
     {\sffamily\bfseries\small\color{white}\colorbox{Bleu}{\strut\ MÉTHODE~\thetcbcounter\ }}%
     \ \ {\sffamily\bfseries\color{Bleu}#2}\par\smallskip},#1}

% ---------- Repères pédagogiques : un filet de couleur, un mot-clé ----------
\newcommand{\definirbloc}[3]{% nom, couleur, libellé
  \newtcolorbox{#1}[1][]{enhanced,breakable,blanker,
    borderline west={1.4pt}{0pt}{#2},left=10pt,right=2pt,top=1pt,bottom=1pt,
    before skip=1em,after skip=1em,
    before upper={{\sffamily\bfseries\footnotesize\color{#2}\textls[120]{\MakeUppercase{#3}}}\quad},##1}}
\definirbloc{reflexe}{Bleu}{Réflexe}
\definirbloc{astuce}{Ocre}{Astuce}
\definirbloc{attention}{Brique}{Attention}
\definirbloc{enpratique}{BleuMoyen}{En pratique}
\definirbloc{plusloin}{Gris}{Pour aller plus loin}
\definirbloc{strategie}{Bleu}{Stratégie}
\definirbloc{lecture}{BleuMoyen}{Lire le schéma}
% compatibilité avec les parties reprises
\newenvironment{garde}{\begin{attention}}{\end{attention}}
\newenvironment{complementHP}{\begin{plusloin}}{\end{plusloin}}
\newenvironment{interpretation}{\begin{enpratique}}{\end{enpratique}}

% ---------- Exercices, corrigés ----------
\newcounter{exo}[section]
\renewcommand{\theexo}{\thesection.\arabic{exo}}
\newcommand{\niveau}[1]{{\color{Ocre}\footnotesize
  \ifcase#1\or$\bullet\circ\circ\circ\circ$\or$\bullet\bullet\circ\circ\circ$\or$\bullet\bullet\bullet\circ\circ$%
  \or$\bullet\bullet\bullet\bullet\circ$\or$\bullet\bullet\bullet\bullet\bullet$\fi}}
\newcommand{\serie}[1]{\par\needspace{8\baselineskip}\vspace{1.5em}%
  \noindent{\sffamily\bfseries\color{Bleu}#1}\par\nobreak\vspace{-0.4em}%
  \noindent{\color{Filet}\rule{\linewidth}{0.4pt}}\par\nobreak\vspace{0.1em}}
\newenvironment{exo}[3]{%
  \par\needspace{5\baselineskip}\vspace{0.75em}\refstepcounter{exo}\label{e:#1}%
  \def\exolab{#1}%
  \noindent{\sffamily\bfseries\color{Bleu}Exercice~\theexo}\quad\niveau{#2}%
  \hfill{\footnotesize\sffamily\color{Gris}#3}\par\nobreak\vspace{0.15em}}%
  {\par\nopagebreak\noindent\hfill{\scriptsize\sffamily\hyperref[c:\exolab]{\color{Gris}corrigé p.~\pageref{c:\exolab}}}\par}
\newenvironment{corrige}[1]{%
  \par\needspace{6\baselineskip}\vspace{1em}\phantomsection\label{c:#1}%
  \noindent{\sffamily\bfseries\color{Bleu}Exercice~\ref{e:#1}}\quad
  {\scriptsize\sffamily\hyperref[e:#1]{\color{Gris}énoncé p.~\pageref{e:#1}}}\par\nobreak\vspace{0.2em}}%
  {\hfill\textcolor{BleuClair}{$\square$}\par}
\newcommand{\etape}[1]{\par\smallskip\noindent{\sffamily\bfseries\small\color{BleuMoyen}#1.}\ }
\newenvironment{solution}{\par\noindent{\sffamily\bfseries\small\color{BleuMoyen}Solution.}\ }{\hfill\textcolor{BleuClair}{$\square$}\par}
\newcommand{\chercher}{\par\smallskip\noindent{\color{Filet}\rule[0.5ex]{1.2cm}{0.4pt}}\
  {\footnotesize\sffamily\itshape\color{Gris}laisser chercher l'étudiant avant de poursuivre}\par\smallskip}

% ---------- Guide de séance ----------

\newcommand{\marge}[1]{\leavevmode\marginpar{\raggedleft\scriptsize\sffamily\color{Gris}\textls[100]{\MakeUppercase{#1}}}\ignorespaces}
\newcommand{\Etape}[4]{%
  \subsection[Étape #1 --- #3 {\color{Gris}(#2)}]{#3}\label{etape:#1}%
  \noindent{\small\sffamily\color{Gris}{\color{Ocre}\bfseries Étape #1}\quad·\quad #2\quad·\quad #4}\par\medskip}
\newcommand{\question}[1]{\par\smallskip\noindent{\sffamily\bfseries\small\color{BleuMoyen}Question.}\ \textit{#1}\par\smallskip}
\newcommand{\rappel}[1]{\par\smallskip\noindent\phantomsection\addcontentsline{rap}{rapentry}{#1}%
  {\sffamily\bfseries\scriptsize\colorbox{Fond}{\color{Gris}\strut RAPPEL}}\ {\sffamily\small\color{Gris}#1.}\ }
\newcommand{\pourquoi}{\par\noindent{\sffamily\bfseries\small\color{Ocre}Pourquoi ?}\ }

% Ouverture de chapitre : \ouverture{objectifs}{prérequis}{durée}
\newcommand{\ouverture}[3]{%
  \begin{tcolorbox}[enhanced,colback=Fond,colframe=Fond,arc=0pt,boxrule=0pt,left=12pt,right=12pt,top=10pt,bottom=10pt,
     before skip=1.2em,after skip=1.6em]
  \begin{minipage}[t]{0.56\linewidth}
  {\sffamily\bfseries\small\color{Bleu}OBJECTIFS}\par\vspace{0.3em}\small #1
  \end{minipage}\hfill
  \begin{minipage}[t]{0.38\linewidth}
  {\sffamily\bfseries\small\color{Bleu}PRÉREQUIS}\par\vspace{0.3em}\small #2
  \ifx\relax#3\relax\else\par\vspace{0.6em}{\sffamily\bfseries\small\color{Bleu}DURÉE}\par\vspace{0.2em}\small #3\fi
  \end{minipage}
  \end{tcolorbox}}
% Fin de chapitre
\newenvironment{essentiel}{\par\needspace{12\baselineskip}\vspace{1.5em}%
  \begin{tcolorbox}[enhanced,unbreakable,colback=white,colframe=white,boxrule=0pt,arc=0pt,
   borderline north={1.4pt}{0pt}{Ocre},left=2pt,right=2pt,top=8pt,bottom=4pt]
  {\sffamily\bfseries\color{Ocre}L'ESSENTIEL DU CHAPITRE}\par\medskip}{\end{tcolorbox}}

% ---------- Notations ----------
\newcommand{\K}{\mathbb{K}}
\newcommand{\R}{\mathbb{R}}
\newcommand{\C}{\mathbb{C}}
\newcommand{\N}{\mathbb{N}}
\newcommand{\Z}{\mathbb{Z}}
\newcommand{\Mn}[1]{\mathcal{M}_{n}(#1)}
\newcommand{\GLn}[1]{\mathrm{GL}_n(#1)}
\newcommand{\LE}{\mathcal{L}(E)}
\newcommand{\rg}{\operatorname{rg}}
\newcommand{\vect}{\operatorname{Vect}}
\newcommand{\Mat}{\operatorname{Mat}}
\newcommand{\Id}{\operatorname{Id}}
\newcommand{\sh}{\operatorname{sh}}
\newcommand{\ch}{\operatorname{ch}}
\renewcommand{\th}{\operatorname{th}}
\DeclareMathOperator{\argsh}{argsh}
\newcommand{\vv}[1]{\overrightarrow{#1}}
\newcommand{\ve}[1]{\vec{#1}}
\newcommand{\ient}[2]{\llbracket #1,#2\rrbracket}
\newcommand{\dt}[1]{\left|\begin{matrix}#1\end{matrix}\right|}
\newcommand{\dis}{\displaystyle}

% ---------- Figures ----------
\pgfplotsset{trig format plots=rad}
\pgfplotsset{courbe/.style={axis lines=middle,axis equal image,
  axis line style={Gris,-{Stealth[length=5pt]}},
  xlabel={$x$},ylabel={$y$},ticklabel style={font=\scriptsize\sffamily,color=Gris},
  every axis x label/.style={at={(ticklabel* cs:1)},anchor=west,font=\small},
  every axis y label/.style={at={(ticklabel* cs:1)},anchor=south,font=\small},
  samples=300,smooth,width=0.6\linewidth}}
\tikzset{arcetude/.style={Bleu,line width=1.25pt},
         reste/.style={BleuClair,line width=1pt},
         tgte/.style={Ocre,line width=0.9pt},
         asy/.style={Gris,dashed,line width=0.6pt},
         pt/.style={circle,fill=Ink,inner sep=1.2pt}}

% ---------- Navigation ----------
\makeatletter
\newcommand*\l@rapentry{\@dottedtocline{1}{0em}{0em}}
\newcommand*\l@metentry{\@dottedtocline{1}{0em}{0em}}
\newcommand{\listemethodes}{\@starttoc{met}}
\newcommand{\listerappels}{\@starttoc{rap}}
\makeatother
% onglets de chapitre sur le bord droit
\usepackage{eso-pic}
\newif\ifonglets \ongletsfalse
\newcommand{\dessineronglet}{%
  \edef\ongcol{\ongletcouleur}%
  \edef\ongnum{\grandnum}%
  \pgfmathsetmacro{\ongy}{3.2+2.3*(\ongnum-1)}%
  \begin{tikzpicture}[remember picture,overlay]
  \fill[\ongcol] ([xshift=-0.75cm,yshift=-\ongy cm]current page.north east)
     rectangle ([yshift=-\ongy cm-1.9cm]current page.north east);
  \node[text=white,font=\sffamily\bfseries\large] at ([xshift=-0.37cm,yshift=-\ongy cm-0.95cm]current page.north east)
     {\ongnum};
  \end{tikzpicture}}
\AddToShipoutPictureBG{\ifonglets\ifnum\value{section}>0 \dessineronglet\fi\fi}
% blocs de navigation et d'enseignement
\definirbloc{commencer}{Gris}{Si vous commencez ici}
\definirbloc{horsprogramme}{Brique}{Hors programme}
\newcommand{\HP}{\,{\sffamily\bfseries\scriptsize\textcolor{Brique}{[HP]}}}
\newcommand{\source}[1]{{\footnotesize\sffamily\color{Gris}\,(#1)}}
\newcommand{\coupdepouce}[1]{\par\vspace{-0.2em}\noindent{\footnotesize\sffamily\color{Gris}%
  \textbf{Coup de pouce} (à donner si l'étudiant bloque) : #1}\par}
\newcommand{\pointarret}[1]{\par\vspace{1em}\noindent{\color{Filet}\rule[0.5ex]{\linewidth}{0.4pt}}\par\nobreak
  \noindent{\footnotesize\sffamily\color{Ocre}\bfseries POINT D'ARRÊT POSSIBLE}\quad{\footnotesize\sffamily\color{Gris}#1}\par\vspace{0.6em}}
\newcommand{\demarche}[1]{\noindent{\sffamily\bfseries\small\color{BleuMoyen}Démarche.}\ \textit{#1}\par\smallskip}

% ---------- Corrigés : idées, astuces, pièges, à retenir ----------
\newcommand{\choix}[1]{\par\noindent{\sffamily\bfseries\small\color{BleuMoyen}Choix de la méthode.}\ #1\par\smallskip}
\newcommand{\gain}[1]{\par\noindent{\sffamily\bfseries\small\color{Ocre}Astuce.}\ #1\par}
\newcommand{\piege}[1]{\par\noindent{\sffamily\bfseries\small\color{Brique}Piège.}\ #1\par}
\newtcolorbox{retenirbox}{enhanced,breakable,colback=Fond,colframe=Fond,arc=0pt,boxrule=0pt,
  borderline west={2pt}{0pt}{Ocre},left=8pt,right=6pt,top=4pt,bottom=4pt,before skip=0.7em,after skip=0.5em}
\newcommand{\retenir}[1]{\begin{retenirbox}{\sffamily\bfseries\footnotesize\color{Ocre}À RETENIR}\quad #1\end{retenirbox}}
\newcommand{\TD}[1]{{\color{Ocre}\bfseries TD\,#1}\ ·\ }
\newcommand{\q}[1]{\par\smallskip\noindent{\sffamily\bfseries\color{Bleu}#1}\ }
\providecommand{\iddots}{\mathinner{\mkern1mu\raise1pt\hbox{.}\mkern2mu\raise4pt\hbox{.}\mkern2mu\raise7pt\hbox{.}\mkern1mu}}


\fancyfoot[L]{}
\newcommand{\Partie}[1]{\par\needspace{5\baselineskip}\vspace{1.1em}\noindent
  {\sffamily\bfseries\large\color{Bleu}#1}\par\vspace{2pt}{\color{Filet}\rule{\linewidth}{0.4pt}}\par\smallskip}
\newcommand{\Q}[1]{\par\needspace{3\baselineskip}\medskip\noindent{\sffamily\bfseries\color{Bleu}#1}\enspace}
\setlist[enumerate,1]{label=\arabic*.,leftmargin=1.6em,itemsep=3pt,topsep=3pt}
\newcommand{\bandeau}[2]{%
\begin{tikzpicture}[remember picture,overlay]
\fill[Bleu] (current page.north west) rectangle ([yshift=-4.2cm]current page.north east);
\node[anchor=north west,text=white,font=\sffamily\footnotesize,inner sep=0] at ([xshift=2.2cm,yshift=-1.1cm]current page.north west)
  {\textls[120]{MATHÉMATIQUES · CPGE PT · ENTRAÎNEMENT}};
\node[anchor=north west,text=white,font=\sffamily\bfseries\LARGE,inner sep=0] at ([xshift=2.2cm,yshift=-1.75cm]current page.north west)
  {Séries numériques et séries entières};
\node[anchor=north west,text=white,font=\sffamily\normalsize,inner sep=0] at ([xshift=2.2cm,yshift=-2.75cm]current page.north west)
  {#1};
\node[anchor=north east,text=white,font=\sffamily\scshape\large,inner sep=0] at ([xshift=-2.2cm,yshift=-1.05cm]current page.north east)
  {Le Collé};
\node[anchor=north east,text=white,font=\sffamily\small,inner sep=0] at ([xshift=-2.2cm,yshift=-2.8cm]current page.north east)
  {#2};
\end{tikzpicture}
\vspace*{2.9cm}}
\fancyhead[L]{\footnotesize\sffamily\color{Gris}Entraînement · séries numériques et séries entières · corrigé Le Collé}
\begin{document}
\pagestyle{fancy}\thispagestyle{plain}
\bandeau{Samedi 28 novembre 2026 · préparation au DS du lundi 30 novembre}{Corrigé}

\begin{tcolorbox}[enhanced,colback=Fond,colframe=Fond,arc=0pt,boxrule=0pt,left=10pt,right=10pt,top=8pt,bottom=8pt]
\small Corrigé rédigé par Le Collé, au niveau de rédaction attendu en devoir. Tous les résultats ont été vérifiés par calcul formel.
\end{tcolorbox}

\Partie{Exercice 1 — Séries numériques}

\Q{1.(a)} $\frac1{\sqrt n}\to0$ et $1-\cos u\sim\frac{u^2}2$ quand $u\to0$ : $u_n\sim\frac1{2n}$. Les termes sont positifs et $\sum\frac1n$ diverge :
par comparaison de séries à termes positifs, \textbf{$\sum u_n$ diverge}.

\Q{1.(b)} $u_n>0$ et
\[\frac{u_{n+1}}{u_n}=\frac{(n+1)^2\cdot3}{(2n+2)(2n+1)}=\frac{3(n+1)}{2(2n+1)}\xrightarrow[n\to+\infty]{}\frac34<1.\]
D'après la règle de d'Alembert, \textbf{$\sum u_n$ converge}.

\Q{1.(c)} $u_n>0$ et $\ln u_n=(\ln n)^2-n\ln(\ln n)$. Alors
\[\ln\big(n^2u_n\big)=2\ln n+(\ln n)^2-n\ln(\ln n)=-n\ln(\ln n)\left(1-\frac{2\ln n+(\ln n)^2}{n\ln(\ln n)}\right)\xrightarrow[n\to+\infty]{}-\infty,\]
car $\frac{(\ln n)^2}n\to0$ (croissances comparées) et $\ln(\ln n)\to+\infty$. Donc $n^2u_n\to0$, $u_n=o\left(\frac1{n^2}\right)$, et
\textbf{$\sum u_n$ converge} (comparaison à une série de Riemann convergente).

\Q{1.(d)} $u_n=(-1)^na_n$ avec $a_n=\ln\left(1+\frac1n\right)$ positive, décroissante (car $x\mapsto\ln(1+x)$ est croissante et $\frac1n$ décroît) et de
limite nulle. D'après le critère des séries alternées, \textbf{$\sum u_n$ converge}. Elle ne converge pas absolument : $a_n\sim\frac1n$.

\Q{2.(a)} $|v_n|=\frac1{\sqrt n+(-1)^n}$ n'est pas décroissante : $|v_{2p}|=\frac1{\sqrt{2p}+1}<\frac1{\sqrt{2p+1}-1}=|v_{2p+1}|$. Le critère des
séries alternées ne s'applique pas.

\Q{2.(b)} On factorise : $v_n=\frac{(-1)^n}{\sqrt n}\cdot\frac1{1+\frac{(-1)^n}{\sqrt n}}$ et $\frac1{1+u}=1-u+O(u^2)$ quand $u\to0$ :
\[v_n=\frac{(-1)^n}{\sqrt n}\left(1-\frac{(-1)^n}{\sqrt n}+O\left(\frac1n\right)\right)=\frac{(-1)^n}{\sqrt n}-\frac1n+O\left(\frac1{n^{3/2}}\right).\]
$\sum\frac{(-1)^n}{\sqrt n}$ converge (critère des séries alternées), $\sum O\big(n^{-3/2}\big)$ converge absolument (Riemann, $\frac32>1$), et
$\sum\frac1n$ diverge. Si $\sum v_n$ convergeait, $\sum\frac1n=\sum\left(\frac{(-1)^n}{\sqrt n}+O(n^{-3/2})-v_n\right)$ convergerait :
\textbf{$\sum v_n$ diverge}.

\Q{3.} $(n+1)^2-n^2=2n+1$, donc $\frac{2n+1}{n^2(n+1)^2}=\frac1{n^2}-\frac1{(n+1)^2}$. Série télescopique :
\[\sum_{k=1}^n\frac{2k+1}{k^2(k+1)^2}=1-\frac1{(n+1)^2}\xrightarrow[n\to+\infty]{}1.\]
La série converge et \textbf{sa somme vaut $1$}.

\Partie{Exercice 2 — Séries entières}

\Q{1.(a)} Pour $z\neq0$ : $\frac{\ln(n+1)|z|^{n+1}}{\ln(n)|z|^n}=\frac{\ln n+\ln(1+\frac1n)}{\ln n}|z|\to|z|$. D'après la règle de d'Alembert, la série
converge absolument si $|z|<1$ et diverge si $|z|>1$ : \textbf{$R=1$}.

\Q{1.(b)} $\left(1+\frac1n\right)^{n^2}=\exp\left(n^2\ln\left(1+\frac1n\right)\right)=\exp\left(n-\frac12+o(1)\right)$, donc
$a_n\sim e^{-1/2}e^n$. Les séries $\sum a_nz^n$ et $\sum(ez)^n$ ont le même rayon : \textbf{$R=\frac1e$}.

\Q{1.(c)} Série lacunaire. Pour $z\neq0$ : $\frac{|z|^{2n+2}/(2^{n+1}+n+1)}{|z|^{2n}/(2^n+n)}=\frac{2^n+n}{2^{n+1}+n+1}|z|^2\to\frac{|z|^2}2$.
Convergence absolue si $|z|<\sqrt2$, divergence si $|z|>\sqrt2$ : \textbf{$R=\sqrt2$}.

\Q{1.(d)} $\frac{a_{n+1}}{a_n}=\frac{(n+1)!\,(2n)!}{n!\,(2n+2)!}=\frac1{2(2n+1)}\to0$ : \textbf{$R=+\infty$}.

\Q{2.(a)} Pour $x\neq0$ : $\frac{n(2n-1)}{(n+1)(2n+1)}x^2\to x^2$. D'après la règle de d'Alembert, $R=1$. En $x=\pm1$, $\left|\frac{(-1)^{n+1}}{n(2n-1)}\right|\sim\frac1{2n^2}$ :
la série converge absolument. \textbf{$f$ est définie sur $[-1,1]$}.

\Q{2.(b)} $\frac1{n(2n-1)}=\frac2{2n-1}-\frac1n$. Pour $|x|<1$, les deux séries suivantes convergent :
\[\sum_{n\geq1}\frac{(-1)^{n+1}}nx^{2n}=\ln(1+x^2),\qquad\sum_{n\geq1}\frac{(-1)^{n+1}}{2n-1}x^{2n}=x\sum_{k\geq0}\frac{(-1)^k}{2k+1}x^{2k+1}=x\arctan x\]
(DSE usuels, avec $k=n-1$). Par linéarité : \textbf{$f(x)=2x\arctan x-\ln(1+x^2)$} pour $x\in\,]-1,1[$.

\Q{2.(c)} La série converge en $x=1$ ; la fonction somme est continue en $1$ (continuité au bord) :
\[\sum_{n=1}^{+\infty}\frac{(-1)^{n+1}}{n(2n-1)}=f(1)=\lim_{x\to1^-}\big(2x\arctan x-\ln(1+x^2)\big)=\frac\pi2-\ln2.\]

\Q{3.} Pour $|x|<1$, $1+x+x^2=\frac{1-x^3}{1-x}>0$ et $g(x)=\ln(1-x^3)-\ln(1-x)=-\sum_{n\geq1}\frac{x^{3n}}n+\sum_{n\geq1}\frac{x^n}n$. Donc
\[g(x)=\sum_{n\geq1}a_nx^n\qquad\text{avec}\qquad a_n=\frac1n\ \text{si }3\nmid n,\qquad a_n=\frac1n-\frac3n=-\frac2n\ \text{si }3\mid n.\]
$\frac1n\leq|a_n|\leq\frac2n$ : par comparaison avec $\sum\frac{x^n}n$ (rayon $1$), $R\geq1$ et $R\leq1$. \textbf{$R=1$}. Enfin $g^{(9)}(0)=9!\,a_9=9!\cdot\left(-\frac29\right)=-2\cdot8!=-80\,640$.

\Q{4.} Pour $|x|<1$, en dérivant deux fois $\frac1{1-x}=\sum x^n$ : $\sum_{n\geq1}nx^n=\frac x{(1-x)^2}$ et $\sum_{n\geq0}n^2x^n=\frac{x(1+x)}{(1-x)^3}$
(avec $n^2=(n+1)(n+2)-3(n+1)+1$). En $x=\frac12$ : $\displaystyle\sum_{n=1}^{+\infty}\frac{n^2}{2^n}=\frac{\frac12\cdot\frac32}{\frac18}=6$.

\Partie{Exercice 3 — Une équation différentielle}

\Q{1.(a)} Sur $]-R,R[$, $y$ est de classe $\mathcal C^\infty$ et se dérive terme à terme :
\[xy''+2y'=\sum_{n\geq2}n(n-1)a_nx^{n-1}+\sum_{n\geq1}2na_nx^{n-1}=\sum_{n\geq1}n(n+1)a_nx^{n-1}=\sum_{n\geq0}(n+1)(n+2)a_{n+1}x^n,\]
et $xy=\sum_{n\geq1}a_{n-1}x^n$. Donc $y$ est solution sur $]-R,R[$ si et seulement si
\[2a_1+\sum_{n\geq1}\big((n+1)(n+2)a_{n+1}+a_{n-1}\big)x^n=0.\]
Par unicité du développement en série entière : $a_1=0$ et $(n+1)(n+2)a_{n+1}+a_{n-1}=0$ pour tout $n\geq1$.

\Q{1.(b)} De $a_1=0$ et $a_{n+1}=-\frac{a_{n-1}}{(n+1)(n+2)}$, tous les coefficients impairs sont nuls. Pour $n=2p-1$ : $a_{2p}=-\frac{a_{2p-2}}{2p(2p+1)}$,
donc par récurrence $a_{2p}=\frac{(-1)^pa_0}{(2p+1)!}$. Pour $x\neq0$, la règle de d'Alembert sur le terme présent donne
$\frac{|x|^2}{(2p+2)(2p+3)}\to0$ : \textbf{$R=+\infty$}. Réciproquement, cette série vérifie la relation de récurrence, donc sa somme est solution
sur $\R$. Pour $x\neq0$ :
\[y(x)=a_0\sum_{p\geq0}\frac{(-1)^px^{2p}}{(2p+1)!}=a_0\,\frac{\sin x}x,\]
et $y(0)=a_0$. Les solutions DSE sont les multiples de la fonction $\varphi$ égale à $\frac{\sin x}x$ pour $x\neq0$ et à $1$ en $0$.

\Q{2.} Sur $I$, $y=\frac zx$ avec $z$ deux fois dérivable si et seulement si $y$ l'est. $z'=y+xy'$ et $z''=2y'+xy''$, donc
$xy''+2y'+xy=z''+z$. Ainsi $y$ est solution de $(E)$ sur $I$ si et seulement si $z''+z=0$, c'est-à-dire $z(x)=A\cos x+B\sin x$. Les solutions
sur $]0,+\infty[$ sont
\[y(x)=\frac{A\cos x+B\sin x}x,\qquad(A,B)\in\R^2,\]
et le même calcul sur $]-\infty,0[$ donne $y(x)=\frac{C\cos x+D\sin x}x$, $(C,D)\in\R^2$.

\Q{3.} Soit $y$ une solution de classe $\mathcal C^2$ sur $\R$. Ses restrictions à $]0,+\infty[$ et $]-\infty,0[$ sont de la forme précédente.
$y$ est continue en $0$, donc bornée au voisinage de $0$. Or $\frac{A\cos x+B\sin x}x\sim\frac Ax$ quand $x\to0^+$ si $A\neq0$ : donc $A=0$, et de
même $C=0$. Alors $y(x)\to B$ en $0^+$ et $y(x)\to D$ en $0^-$, donc $B=D=y(0)$ : $y=B\varphi$.

Réciproquement, $\varphi$ est la somme d'une série entière de rayon infini : elle est $\mathcal C^\infty$ sur $\R$ et solution de $(E)$ (question 1).
Les solutions de classe $\mathcal C^2$ sur $\R$ sont donc les $B\varphi$, $B\in\R$ : elles forment une \textbf{droite} vectorielle, et non un espace
de dimension $2$.

Le théorème de Cauchy s'applique à une équation $y''=a(x)y'+b(x)y$ à coefficients continus sur l'intervalle. Ici le coefficient de $y''$
est $x$, qui s'annule en $0$ : on ne peut pas écrire $(E)$ sous forme résolue sur $\R$, ni sur aucun intervalle contenant $0$.
\end{document}
